\subsection{Definition of factorial domains}

DEFINITION 1. A Krull domain all of whose divisorial ideals are principal is called a factorial (or unique factorization) domain.

In other words, the group of divisor classes (§ 1, no. 2) is reduced to 0.

\textbf{Examples}\par

\begin{enumerate}
\def\labelenumi{(\arabic{enumi})}
\item
  Every principal ideal domain is factorial (and, recall, is a Dedekind domain). Conversely, every factorial Dedekind domain is a principal ideal domain by § 2, no. 2, Theorem 1 (c).
\item
  In particular, if K is a field, the rings K{[}X{]} and K{[}{[}X{]}{]} are factorial domains (see Theorem 2 and Proposition 8 below for generalizations).
\item
  * The local ring of a simple point of an algebraic variety is a factorial domain. The ring of germs of functions analytic at the origin of \(\mathbf{C}^n\) is a factorial domain. *
\end{enumerate}

